% !TeX root = ../thesis.tex
% Kapitel 4
\chapter{Methodik und Umsetzung}
\label{cha:methodik}

\section{Datenzugang und Datenschutz}
\label{sec:datenschutz}
\todo{Read-only-Extraktion, Anonymisierung bei der Extraktion, Kontrollen aus \code{riskControlMapping}, nur Aggregate in der Arbeit.}

\section{Population, Vorfilter und Zielgröße}
\label{sec:daten}
\todo{Aktive Risiken mit mindestens einer Familie; Vorfilter ohne Zielgröße; Familienbildung; Kennzahlen aus \code{run\_meta.json}.}

\section{Pipeline und Evaluationsprotokoll}
\label{sec:protokoll}
\begin{figure}[htbp]
  \centering
  \begin{tikzpicture}[
      node distance=0.35cm and 0.35cm,
      box/.style={draw, rounded corners, align=center, minimum height=0.9cm,
                  text width=1.6cm, font=\scriptsize},
      inner/.style={box, fill=gray!12},
      arr/.style={-{Stealth}, thick}]
    \node[box] (load) {load};
    \node[box, right=of load] (pre) {prefilter};
    \node[box, right=of pre]  (lab) {labels};
    \node[box, right=of lab]  (spl) {split};
    \node[box, below=0.9cm of spl] (exp) {experiment\\(5 Durchgänge)};
    \node[box, left=of exp]   (rep) {report};
    \draw[arr] (load) -- (pre);
    \draw[arr] (pre) -- (lab);
    \draw[arr] (lab) -- (spl);
    \draw[arr] (spl) -- (exp);
    \draw[arr] (exp) -- (rep);
    \node[inner, below=1.1cm of rep, xshift=-3.0cm] (sel) {selection};
    \node[inner, right=of sel] (fit) {fit};
    \node[inner, right=of fit] (msk) {draft\_\\simulation};
    \node[inner, right=of msk] (rec) {recommend};
    \node[inner, right=of rec] (eva) {evaluate};
    \draw[arr] (sel) -- (fit);
    \draw[arr] (fit) -- (msk);
    \draw[arr] (msk) -- (rec);
    \draw[arr] (rec) -- (eva);
    \draw[dashed] (exp.south) -- ++(0,-0.45) -| (eva.north);
  \end{tikzpicture}
  \caption[Aufbau der Pipeline]{Äußere Pipeline (oben, einmal) und innere Pipeline
    (unten, je Durchgang).}
  \label{fig:pipeline}
\end{figure}

\begin{algorithm}[htbp]
\caption[Evaluationsprotokoll]{Evaluationsprotokoll über alle Durchgänge}
\label{alg:protokoll}
\begin{algorithmic}[1]
\REQUIRE vorgefilterte Attribute $X$, Familienmatrix $Y$, Durchgangszuordnung $d$,
  Selektoren $\mathcal{S}$, Methoden $\mathcal{M}$, Maskierungsraten $\mathcal{Q}$, $k$
\ENSURE Ergebniszeilen im Langformat
\FOR{$t \leftarrow 0$ \MYTO 4}
  \STATE $(X^{tr}, Y^{tr}, X^{te}, Y^{te}) \leftarrow$ \MYCALL{TrainTest}{$X$, $Y$, $d$, $t$}
  \FORALL{$s \in \mathcal{S}$}
    \STATE $A_s \leftarrow s.$\MYCALL{select}{$X^{tr}$, $Y^{tr}$} \ \algorithmiccomment{sieht nur das Training}
    \FORALL{$m \in \mathcal{M}$}
      \STATE $\mu \leftarrow$ neues Modell $m$; $\mu.$\MYCALL{fit}{$X^{tr}[A_s]$, $Y^{tr}$}
      \FORALL{$q \in \mathcal{Q}$}
        \STATE $\tilde{X} \leftarrow$ \MYCALL{mask}{$X^{te}$, $q$, $\mathit{seed}+t$}
        \STATE $\hat{R} \leftarrow \mu.$\MYCALL{recommend}{$\tilde{X}[A_s]$, $k$} \ \algorithmiccomment{nie $Y^{te}$}
        \STATE schreibe Zeile $(t, s, m, q, |A_s|,$ \MYCALL{Evaluate}{$\hat{R}$, $Y^{te}$}$)$
      \ENDFOR
    \ENDFOR
  \ENDFOR
\ENDFOR
\end{algorithmic}
\end{algorithm}
\todo{Protokoll begründen: 5 stratifizierte Durchgänge, Leakage-Schutz, Seeds.}
Um die Methoden unter identischen Bedingungen zu vergleichen, ist die Untersuchung als Pipeline mit zwei Ebenen umgesetzt (Abbildung~\ref{fig:pipeline}). 
Die äußere Pipeline läuft einmal und bereitet die Daten auf. Die innere Pipeline wird für jeden Durchgang erneut ausgeführt und enthält alle Schritte, die von der Zielgröße abhängen.
Diese Trennung legt zugleich das \emph{Evaluationsprotokoll} fest, also die vorab bestimmte Vorschrift, welcher Schritt welche Daten sehen darf.

Die äußere Pipeline lädt zunächst die Risiken und ihre Kontrollzuordnungen. Anschließend entfernt der Vorfilter Attribute, die für einen Vergleich ungeeignet sind (Abschnitt~\ref{sec:daten}). 
Danach wird die Zielgröße $Y$ gebildet (Definition~\ref{def:ziel}) und die Population auf Risiken mit mindestens einer Kontrollfamilie beschränkt. 
Schließlich werden die Risiken auf fünf Durchgänge aufgeteilt, die Experimente ausgeführt und die Ergebnisse gespeichert.

Die Aufteilung erfolgt nach dem Prinzip der Kreuzvalidierung: Die Risiken werden in fünf gleich große Teile zerlegt. 
In jedem Durchgang dient ein Teil als Testmenge, die übrigen vier als Trainingsmenge, sodass jedes Risiko genau einmal getestet wird. 
Bei der begrenzten Datenmenge nutzt dieses Vorgehen alle Risiken für die Bewertung, anstatt einen Teil dauerhaft zurückzuhalten. 
Da die Kontrollfamilien sehr ungleich verteilt sind, wird die Aufteilung nach Familien stratifiziert. 
So enthält jeder Teil einen ähnlichen Anteil jeder Familie und auch seltene Familien sind in jedem Durchgang vertreten.

In jedem Durchgang durchläuft die innere Pipeline fünf Schritte (Algorithmus~\ref{alg:protokoll}). Zuerst bestimmt jeder Selektor auf der Trainingsmenge eine Attributmenge. 
Für jede Kombination aus Selektor und Methode wird anschließend ein neues Modell auf der Trainingsmenge trainiert. 
Danach wird die Entwurfsphase simuliert, indem ein Teil der Attributwerte der Testrisiken entfernt wird (Abschnitt~\ref{sec:maskierung}).
Das Modell erzeugt für jedes Testrisiko eine Empfehlungsliste $\hat{R}_r$ der Länge $k$, die abschließend mit den tatsächlichen Familien $Y_r$ verglichen wird.

Zentrale Anforderung an das Protokoll ist die Vermeidung von Datenlecks (\emph{Leakage}). Attributselektion und Training sehen ausschließlich die Trainingsmenge. 
Die Empfehlung erhält nur die Attribute der Testrisiken, nie deren Familien. Diese werden erst bei der Bewertung genutzt. 
Andernfalls würden die Ergebnisse die Güte im realen Einsatz überschätzen. Die Pipeline prüft diese Bedingungen vor jedem Schritt, 
zusätzlich stellt ein automatisierter Test sicher, dass kein Testrisiko in Selektion oder Training gelangt.

Für die Reproduzierbarkeit sind alle Parameter in einer zentralen Konfigurationsdatei abgelegt, und sämtliche Zufallsprozesse verwenden feste Startwerte. 
Jeder Lauf speichert seine Ergebnisse zusammen mit der verwendeten Konfiguration in einem eigenen Verzeichnis.

\subsection{Attributselektion}
\label{sec:selektion}
\todo{Selektoren \code{none} und \code{grande} beschreiben, Stabilität über Durchgänge.}

\subsection{Empfehlungsmethoden}
\label{sec:methoden}
\todo{\code{baseline}, \code{knn\_gower} (inkl.\ Confidence) und \code{clustering} beschreiben, sobald implementiert.}

\subsection{Simulation der Entwurfsphase}
\label{sec:maskierung}
\todo{Test-Attribute zu 0/20/40/60\,\% leeren, Training bleibt unverändert.}
